\section{Superspace model sets as geometric primitives}\label{sec:superspace}

Superspace (or ``higher-dimensional embedding'') methods are a standard geometric language for quasi-periodic structures: a pattern that is aperiodic in a low-dimensional ``physical'' space is represented as a slice (or projection) of a periodic structure in a higher-dimensional ambient space. The cut-and-project construction gives this language a precise mathematical form and is the backbone of model-set theory \cite{Meyer1972,Moody2000,BaakeGrimm2013}.

\subsection{Cut-and-project schemes}\label{sec:superspace:cps}

We fix finite-dimensional real vector spaces
\[
E := \RR^{d} \quad (\text{physical space}),\qquad
H := \RR^{n} \quad (\text{internal space}),
\]
and let $G := E \times H$. Let $\pi_E : G \to E$ and $\pi_H : G \to H$ denote the coordinate projections.

\begin{definition}[Cut-and-project scheme]\label{def:cps}
A \emph{cut-and-project scheme} (CPS) consists of $(E,H,\Gamma)$ where $\Gamma \subset G$ is a lattice such that:
\begin{enumerate}[label=(\roman*),leftmargin=*,itemsep=2pt]
  \item $\pi_E|_{\Gamma}$ is injective;
  \item $\pi_H(\Gamma)$ is dense in $H$.
\end{enumerate}
We write $\Gamma_E := \pi_E(\Gamma)$ and define the \emph{star map} ${}^\star : \Gamma_E \to H$ by $x^\star := \pi_H(\gamma)$ where $\gamma \in \Gamma$ is the unique lattice point with $\pi_E(\gamma)=x$.
\end{definition}

In the quasicrystal/superspace context, one often takes $d=3$ and $d+n=6$ (a ``rank-6'' superspace) to capture quasi-periodic structures with six independent reciprocal vectors. We will keep $d$ and $n$ general but use ``6D'' language when emphasizing the GRG-ETE-motivated corpus that inspired this review.

\subsection{The (3+3) decomposition and superspace densities}\label{sec:superspace:33}

The rank-$6$ case is often presented as a $(3+3)$ decomposition of a six-dimensional Euclidean space:
\[
E^{6}=E_{\parallel}^{3}\oplus E_{\perp}^{3},
\]
where $E_{\parallel}^{3}$ is the physical subspace and $E_{\perp}^{3}$ is the internal subspace. A canonical lattice is the hypercubic lattice $\ZZ^{6}\subset E^{6}$, with orthogonal projections $P_{\parallel}$ and $P_{\perp}$ onto the two summands.

In crystallographic superspace language, a three-dimensional quasi-periodic density $\rho$ can be represented as the restriction of a periodic density $\rho_s$ on superspace to a physical section, schematically $\rho(r)=\rho_s(r,0)$ in suitable coordinates \cite{IUCRSuperspace}. In this review we do not need the full crystallographic machinery; the key point is that the superspace description provides a periodic object on which standard geometric and harmonic tools apply, while the observed pattern arises from a slice/projection determined by protocol parameters (orientation and internal phase).

\subsection{Windows, acceptance domains, and model sets}\label{sec:superspace:windows}

The observational content of a CPS is extracted by selecting a compact subset $W \subset H$ (a \emph{window} or \emph{acceptance domain}) that specifies which lattice points are ``accepted'' by the slice.

\begin{definition}[Model set]\label{def:model-set}
Given a CPS $(E,H,\Gamma)$ and a window $W \subset H$, the associated \emph{model set} is
\[
\Lambda(W) := \{x \in \Gamma_E : x^\star \in W\} \subset E.
\]
If $W$ is compact with nonempty interior and $\partial W$ has Haar measure zero, $\Lambda(W)$ is called a \emph{regular model set}.
\end{definition}

\begin{remark}[Internal phase]\label{rem:internal-phase}
In applications, one often includes an internal phase (or offset) $u\in H$ and considers
\[
\Lambda(W,u):=\{x\in\Gamma_E:\ x^\star\in W+u\}.
\]
Varying $u$ corresponds to translating the window in internal space and is a clean way to parameterize ``observer phase'' at the object layer.
\end{remark}

In crystallographic superspace terminology, $W$ can be described via atomic surfaces and acceptance domains, while the physical pattern $\Lambda(W)$ corresponds to the intersection of the superspace periodic object with the physical subspace followed by projection \cite{Janssen2007,Steurer2012}. The review perspective we adopt is to treat $W$ as part of the \emph{protocol layer}: it models finite-resolution selection and internal-phase sensitivity.

\subsection{Acceptance domains and atomic surfaces}\label{sec:superspace:atomic}

Superspace crystallography distinguishes several closely related concepts that are useful for connecting geometry to readout:
\begin{itemize}[leftmargin=*,itemsep=2pt]
  \item \textbf{Acceptance domains (windows)}: the subset $W\subset H$ specifying which lattice points are accepted. The ``intersection method'' and the ``cut-and-project method'' can be regarded as equivalent formulations of this selection mechanism \cite{IUCRITCVolC}.
  \item \textbf{Atomic surfaces}: higher-dimensional objects periodically arranged in superspace whose intersections with the physical subspace generate physical atoms/vertices; their dimension equals the internal-space dimension \cite{IUCRAtomicSurface}.
\end{itemize}
From the review's perspective, these are not merely descriptive labels: they provide a disciplined vocabulary for how a protocol (window choice and boundary resolution) controls the transition from superspace periodicity to physical quasi-periodicity.

\subsection{Invariant measures and typicality}\label{sec:superspace:measures}

Cut-and-project constructions have a natural dynamical system behind them: translations on the compact Abelian group $G/\Gamma$ (a torus when $G$ is Euclidean). For a fixed CPS, varying the internal ``phase'' corresponds to translating the window in $H$ (or translating the physical slice in $G$). Under mild hypotheses, one obtains unique ergodicity and well-defined pattern frequencies \cite{Moody2000,BaakeGrimm2013}.

From the standpoint of this paper, the key point is that \emph{typicality statements} (e.g.\ long-run frequencies being independent of phase for almost all phases) are available at the object layer and can be used to justify protocol-layer statistical laws. This becomes critical when defining information-theoretic time via refinement: entropy rates are almost-sure quantities under ergodic assumptions.

\begin{proposition}[Unique ergodicity and uniform frequencies (informal review statement)]\label{prop:unique-ergodicity}
For regular model sets arising from a CPS with compact window whose boundary has Haar measure zero, the associated dynamical hull is uniquely ergodic. In particular, patch frequencies exist and are uniform (phase-independent) for almost all internal phases.
\cite{Moody2000,BaakeGrimm2013}
\end{proposition}

\subsubsection{Fourier modules and why superspace is spectrally convenient}

One major reason superspace model sets are attractive as geometric primitives is that they come with an explicit reciprocal-space description. Roughly speaking, the Fourier module of a model set is determined by the dual lattice $\Gamma^{\ast}\subset G^{\ast}$ and by the projections dual to $\pi_E$ and $\pi_H$. This makes it possible to predict (and test) where Bragg peaks may occur and how their intensities depend on the window geometry \cite{Hof1995,BaakeGrimm2013}. In a geometry-first reconstruction context, this means that a significant part of the ``observable content'' of the object layer can be stated as a module-theoretic constraint, independent of interpretation.

\subsubsection{Torus parametrization and the ``maximal equicontinuous factor''}

A useful conceptual bridge between superspace geometry and observation protocols is the torus parametrization viewpoint. The compact group $G/\Gamma$ plays the role of a maximal equicontinuous factor for the translation action on the hull of a model set, and the internal phase $u$ can be identified with a point on this torus. From this angle, ``varying the observer'' corresponds to shifting along a compact parameter space, while the observed pattern in $E$ is a factor determined by the window. This language is particularly convenient when readout is modeled as sampling along a one-dimensional factor: a scan is literally a translation on $G/\Gamma$ projected onto a low-dimensional torus factor.

\subsection{Why 6D matters in practice}\label{sec:superspace:6d}

While CPS and model sets are defined in general dimension, the rank-$6$ case is structurally prominent in the corpus we synthesize, because it provides a minimal higher-dimensional periodic scaffold that can encode rich quasi-periodic 3D patterns while still admitting explicit arithmetic and symbolic codings (e.g.\ via one-dimensional torus factors, Sturmian codings, and canonical Fibonacci/Zeckendorf languages). In this review we will therefore repeatedly refer to ``6D superspace'' as a canonical working model, without restricting the mathematical statements to that dimension unless explicitly said.

\subsection{Rational approximants vs.\ irrational limits}\label{sec:superspace:rational-irrational}

A practical issue in both modeling and observation is that irrational orientations/slopes are never realized with infinite precision. In superspace cut-and-project settings this leads to a natural hierarchy of \emph{rational approximants}: a rationally oriented strip (or rationally related internal/physical decomposition) yields periodic patterns, while the irrational limit yields a quasi-periodic model set. Concrete 6D constructions used in the quasicrystal literature exhibit this dichotomy explicitly, including families where a rational $(1/1)$ strip produces a periodic phase and an irrational $(\tau/1)$ strip (with $\tau=(1+\sqrt5)/2$) produces an icosahedral quasicrystalline phase \cite{PRB46_6091}. This ``rational/irrational = periodic/quasiperiodic'' theme will reappear at the protocol level when we discuss internal phase scans and their rational approximations.

\subsection{Representative icosahedral examples in 6D}\label{sec:superspace:examples}

To make the above discussion less abstract, we mention two representative Phys.\ Rev.\ B constructions that illustrate how the same structure can be described equivalently as a 3D quasi-periodic object or as a 6D periodic object:
\begin{itemize}[leftmargin=*,itemsep=2pt]
  \item A six-dimensional structure model for the icosahedral quasicrystal ${\mathrm{Al}}_{6}{\mathrm{CuLi}}_{3}$ explicitly develops the 6D periodic description and shows how the physical structure is obtained by an appropriate 3D section through the higher-dimensional object \cite{PRB43_929}.
  \item A six-dimensional crystal-structure model for i-(Al--Mn--Si) and $\alpha$-(Al--Mn--Si) describes how changing strip orientations produces either an icosahedral quasicrystal or a periodic approximant, with the golden ratio $\tau$ entering through specific partition ratios along face diagonals and through the irrational strip \cite{PRB46_6091}.
\end{itemize}

For the purposes of a geometry-first reconstruction review, these examples serve as a reality check: superspace cut-and-project is not only an abstract mathematical construction but a standard tool in the analysis of real quasi-periodic materials. This strengthens the case for using superspace model sets as a ``geometric backbone'' on which one can define protocol-level readout and information-theoretic time without immediately committing to a specific microscopic physical ontology.

\subsection{What superspace does (and does not) imply for spacetime geometry}\label{sec:superspace:spacetime}

Because this review is written for a GRG topical collection, it is worth making a careful conceptual clarification. A superspace CPS is \emph{not} a claim that physical spacetime has extra spatial dimensions in the Kaluza--Klein sense, nor that the internal space is physically traversed. Instead, superspace provides a mathematically controlled \emph{representation space} in which a complicated lower-dimensional pattern is expressed as a slice of a periodic object.

From the standpoint of reconstruction, the value of superspace is structural:
\begin{itemize}[leftmargin=*,itemsep=2pt]
  \item It yields a high-symmetry object layer with explicit dual modules and well-controlled ergodic properties.
  \item It builds in a visible/hidden split that can be interpreted operationally: internal degrees of freedom are only accessible through projection and finite-resolution readout.
  \item It provides a natural arena for ``emergence by quotient'': the observed structure is a quotient (projection/cut) of a richer periodic geometry.
\end{itemize}

At the same time, one should avoid importing physical conclusions that are not warranted by the mathematics. In particular:
\begin{itemize}[leftmargin=*,itemsep=2pt]
  \item Superspace periodicity does not by itself imply dynamical laws on the observed space, let alone Einsteinian dynamics.
  \item The internal phase parameter $u$ is a mathematically natural coordinate on a torus factor; interpreting it as an ``observer state'' is a protocol-level modeling choice.
  \item The cut-and-project mechanism produces quasi-periodicity and module constraints, but additional assumptions are needed to connect these constraints to spacetime curvature, causal cones, or quantum superposition.
\end{itemize}

The remainder of this review is written in a way that respects these boundaries: we use superspace to supply an explicit object layer, and we treat readout, stabilization, time, and causality as protocol-defined interfaces whose physical interpretation must be argued separately.

